%% Cross-check of the He I 1.0833 um triplet optical-depth normalization
%% in Draine (2026), against MoCHII and standard atomic physics.
%% Author: Kwang-Il Seon
%% Written for my_memo.cls.
\documentclass[english,a4paper]{my_memo}
\synctex=1
\usepackage{amsmath}
\usepackage{amssymb}
\usepackage{microtype}
\usepackage{booktabs}
\usepackage{array}
\hypersetup{hidelinks}

\newcommand{\code}[1]{\texttt{#1}}
\newcommand{\HeI}{He\,\textsc{i}}
\newcommand{\HII}{H\,\textsc{ii}}
\newcommand{\tSSS}{2\,^3\!S_1}
\newcommand{\tPPP}{2\,^3\!P^o_J}
\newcommand{\tautot}{\tau_{\rm tot}}
\newcommand{\gu}{g_u}
\newcommand{\gl}{g_\ell}
\newcolumntype{C}{>{\centering\arraybackslash}p{0.15\textwidth}}

\begin{document}

\title{The He\,\textsc{i} 1.0833\,$\mu$m Triplet Optical Depth in Draine (2026):\\
A Factor-of-Three Cross-Check against MoCHII}
\author{Kwang-Il Seon}
\date{\docmoddate}
\maketitle

\begin{abstract}
While cross-checking the MoCHII metastable \HeI{} $\tSSS$ diagnostic
against B.\,T.\ Draine's recent study of resonant scattering of the
\HeI{} $1.0833\,\mu$m triplet in \HII{} regions (2026, ApJ 999, 3;
arXiv:2601.11710), I find that the two treatments agree on the
metastable population physics but disagree by exactly a factor of three
in the line-center scattering optical depth.  The paper's optical-depth
normalization (its Eqs.~8, 9, 11, and 12) is three times larger than the
standard atomic-physics result for a given $\tSSS$ column density.  The
discrepancy traces to the upper-level statistical weights listed in the
paper's Table~1, which are $g_u = 3, 9, 15$ where the true fine-structure
degeneracies of $\tPPP$ are $2J+1 = 1, 3, 5$.  The printed weights are
internally inconsistent with the NIST $A$- and $f$-values that the same
table cites: those values are mutually consistent \emph{only} with
$g_u = 2J+1$ and the $\tSSS$ lower-level weight $\gl = 3$.  MoCHII's
$\tau_0$ normalization, the p-winds/Oklop\v{c}i\'c \& Hirata (2018)
exoplanet standard, and the LaRT $1.0833\,\mu$m module all use the
$2J+1$ weights and agree with one another.  Because \HII{} regions are
optically thick to the triplet under either normalization, the paper's
qualitative conclusions (broad multi-peaked profiles, an overall
blueshift, and dust-absorption enhancement by resonant trapping) are
unaffected, but the quantitative $\tautot$ thresholds and the derived
column-density scale should be read a factor of three lower.  All
numerical claims below were recomputed independently.
\end{abstract}

%==========================================================================
\section{Context}
The \HeI{} $1.0833\,\mu$m triplet arises from the $\tSSS$--$\tPPP$
transition and is fed by recombination of He$^+$ into the triplet system,
which cascades into the long-lived metastable $\tSSS$ level.  MoCHII
computes the $\tSSS$ population as a closed-form local balance in every
leaf on the converged state (the Oklop\v{c}i\'c \& Hirata 2018 balance,
hereafter OH18), and derives the $1.0833\,\mu$m line-center opacity from
it; the design and gate results are recorded in
\code{docs/HEI\_10833\_PLAN.md}.  Draine (2026) treats the same level as
the resonant scatterer in a Monte-Carlo transfer calculation of the
$1.0833\,\mu$m emission profile emerging from an \HII{} region.  Since the
two efforts model the same atomic level in overlapping regimes, they form
a natural mutual check.  The population physics agrees; the optical-depth
normalization does not.  This memo records the cross-check.

Throughout, $\tSSS$ is the metastable lower level (statistical weight
$\gl = 2J+1 = 3$), and $\tPPP$ are the three upper fine-structure levels
$J = 0, 1, 2$ (true weights $2J+1 = 1, 3, 5$).  The Doppler broadening
parameter is $b$, quoted in the paper's fiducial units of $10$\,km\,s$^{-1}$,
and $N \equiv N(\tSSS)$ is the center-to-edge column density of metastable
helium.

%==========================================================================
\section{Where the paper and MoCHII agree}
The metastable population physics is common to both treatments.  Draine's
Eq.~(1),
\begin{equation}
  n(\tSSS) = \frac{\alpha_{\rm trip}\, n_e\, n({\rm He}^+)}
                  {A_{ms} + n_e k_d},
\end{equation}
is the OH18 balance MoCHII solves, reduced to its two dominant terms:
recombination into the triplet system as the source, and radiative decay
of $\tSSS$ plus collisional transfer to the singlet system as the sink.
Draine notes that $\tSSS$ photoionization ($h\nu > 4.77$\,eV) and Penning
ionization with H$^0$ are usually subdominant, exactly as MoCHII finds in
the dense \HII{} regime.  The rate coefficients agree to a few percent:
\begin{itemize}\itemsep2pt
  \item $\alpha_{\rm trip}(9000\,{\rm K}) = 2.42\times10^{-13}$\,cm$^3$\,s$^{-1}$
        (Del Zanna \& Storey 2022) vs.\ MoCHII/OH18
        $2.28\times10^{-13}$ at the same temperature --- within 6\%;
  \item $k_d = 3.3\times10^{-8}$\,cm$^3$\,s$^{-1}$ (Bray et al.\ 2000) vs.\
        MoCHII's $q_{31a}+q_{31b} = 3.2\times10^{-8}$ from the same Bray
        source --- within $\sim$3\%;
  \item critical density $n_{\rm crit} = A_{ms}/k_d$: $3855$\,cm$^{-3}$
        (Draine) vs.\ $\sim$$3.9\times10^{3}$\,cm$^{-3}$ (MoCHII, $3945$)
        --- the small difference following directly from $k_d$;
  \item Draine's Eq.~(4) center-to-edge $N(\tSSS)$ formula reproduces the
        MoCHII low-density example (Q$_{48}=10$, $n_e = 110$\,cm$^{-3}$,
        $T_e = 10^4$\,K; $N \sim 1.4\times10^{13}$\,cm$^{-2}$) to within
        $\sim$25\%.
\end{itemize}
On the population side the two codes are consistent within their input
atomic-data spread.

%==========================================================================
\section{The factor-of-three discrepancy}
The disagreement is entirely in the mapping from a metastable column
density $N$ to a scattering optical depth.  Draine's Eqs.~(8) and (9)
give the center-to-edge line-center optical depth for
$\tSSS \to \tPPP$ absorption as
\begin{equation}
  \tau_0(\tSSS \to \tPPP)
    = (2J+1)\,\frac{A_{u\ell}\,\lambda_{u\ell}^3\,N}{8\pi^{3/2} b}
    = 29.2\,(2J+1)
      \left(\frac{N}{10^{14}\,{\rm cm^{-2}}}\right)
      \left(\frac{10\,{\rm km\,s^{-1}}}{b}\right),
  \label{eq:draine8}
\end{equation}
and summing the three components (his Eq.~11) yields
\begin{equation}
  \tautot \equiv \sum_{J=0}^{2}\tau_0(\tSSS \to \tPPP)
    = 262
      \left(\frac{N}{10^{14}\,{\rm cm^{-2}}}\right)
      \left(\frac{10\,{\rm km\,s^{-1}}}{b}\right).
  \label{eq:draine11}
\end{equation}
The standard line-center cross section for a Doppler-broadened line,
written either in oscillator-strength or in Einstein-$A$ form,
\begin{equation}
  \sigma_0
    = \frac{\sqrt{\pi}\,e^2}{m_e c}\, f_{\ell u}\,\frac{\lambda_{u\ell}}{b}
    = \frac{\gu}{\gl}\,\frac{A_{u\ell}\,\lambda_{u\ell}^3}{8\pi^{3/2} b},
  \label{eq:sigma0}
\end{equation}
carries the ratio $\gu/\gl$.  With the true weights $\gu = 2J+1$ and
$\gl = 3$, and the Table-1 wavelengths and $A$-values, the three
components sum to
\begin{equation}
  \tautot^{\rm std}
    = 87.5
      \left(\frac{N}{10^{14}\,{\rm cm^{-2}}}\right)
      \left(\frac{10\,{\rm km\,s^{-1}}}{b}\right),
\end{equation}
which I confirmed to be identical whether computed from the
oscillator-strength form or the $A$-value form of Eq.~(\ref{eq:sigma0}).
Component by component (per $10^{14}$\,cm$^{-2}$, $b = 10$\,km\,s$^{-1}$):

\begin{center}
\begin{tabular}{cccc}
\toprule
$J$ & standard $\tau_0$ ($\gu = 2J{+}1$, $\gl = 3$)
    & paper's Eq.~(8)
    & ratio \\
\midrule
0 & \phantom{0}9.72 & \phantom{0}29.2 & 3.00 \\
1 & 29.16           & \phantom{0}87.5 & 3.00 \\
2 & 48.60           & 145.8           & 3.00 \\
\midrule
sum & 87.5 & 262.4 & 3.00 \\
\bottomrule
\end{tabular}
\end{center}

The ratio is $262.4/87.5 = 3.0000$ to the precision of the input data.
The paper's Eq.~(8) is the standard $\sigma_0 N$ of
Eq.~(\ref{eq:sigma0}) with $\gu \to (2J+1)$ but with the lower-level
weight $\gl$ set to unity rather than to its true value $\gl = 3$; the
missing denominator is the source of the factor of three.

%==========================================================================
\section{Root cause: the statistical weights}
The missing $\gl$ is visible directly in the paper's own Table~1, whose
$g_u$ column lists $3, 9, 15$ where the fine-structure levels $\tPPP$ have
degeneracies $2J+1 = 1, 3, 5$.  Table~\ref{tab:weights} places the two
side by side.

\begin{table}[h]
\centering
\caption{The \HeI{} $1.0833\,\mu$m triplet: statistical weights and
oscillator strengths.  The paper's Table~1 lists $g_u = 3, 9, 15$, which
are $3\times$ the true fine-structure degeneracies $2J+1$.  The $f_{\ell u}$
values, computed from the tabulated $A_{u\ell}$ with the true
$\gu = 2J+1$ and $\gl = 3$, reproduce the NIST oscillator strengths (their
1\,:\,3\,:\,5 ratio and sum $0.539$) exactly.}
\label{tab:weights}
\begin{tabular}{ccccccc}
\toprule
$u$ & $J$ & Draine $g_u$ & true $2J+1$
    & $A_{u\ell}$ (s$^{-1}$) & $\lambda_{\rm vac}$ ($\mu$m)
    & $f_{\ell u}$ (this work $=$ NIST) \\
\midrule
$2\,^3\!P^o_0$ & 0 & \phantom{0}3 & 1 & $1.0216\times10^7$ & 1.083206 & 0.05990 \\
$2\,^3\!P^o_1$ & 1 & \phantom{0}9 & 3 & $1.0216\times10^7$ & 1.083322 & 0.17974 \\
$2\,^3\!P^o_2$ & 2 & 15          & 5 & $1.0216\times10^7$ & 1.083331 & 0.29958 \\
\midrule
sum & & 27 & 9 & & & 0.539 \\
\bottomrule
\end{tabular}
\end{table}

The most likely origin of $3, 9, 15$ is a redundant multiplication by the
spin multiplicity $2S+1 = 3$: the fine-structure weight $2J+1$ already
partitions the full term degeneracy $(2S+1)(2L+1) = 9$ among the three
$J$ levels ($\sum_J (2J+1) = 9$), so multiplying each $2J+1$ again by
$2S+1 = 3$ double-counts the triplet spin and yields
$3, 9, 15$ (sum $27 = 3\times9$).  In the paper's Eq.~(8) this
same factor of three appears as the omitted $\gl = 3$: using the standard
$\sigma_0 = (\gu/\gl)\cdots$ with the inflated $\gu = 3(2J+1)$ and the
correct $\gl = 3$ gives $\gu/\gl = (2J+1)$, exactly the bare $(2J+1)$
printed in the paper's Eq.~(8).  The two readings --- ``Table~1 $g_u$
inflated by 3'' and ``$\gl$ omitted from Eq.~8'' --- are the same
statement.

%==========================================================================
\section{Proof by NIST self-consistency}
Independently of any external code, the printed weights are inconsistent
with the $A$- and $f$-values Table~1 itself cites (NIST ASD; Kramida et
al.\ 2024).  The exact relation between an Einstein $A$ and an absorption
oscillator strength is
\begin{equation}
  \gl\, f_{\ell u}
    = \frac{m_e c}{8\pi^2 e^2}\,\lambda_{u\ell}^2\, \gu\, A_{u\ell}
  \quad\Longrightarrow\quad
  f_{\ell u} = \frac{\gu}{\gl}\,
    \frac{m_e c\,\lambda_{u\ell}^2 A_{u\ell}}{8\pi^2 e^2}.
  \label{eq:Af}
\end{equation}
Substituting the Table-1 $A_{u\ell} = 1.0216\times10^7$\,s$^{-1}$ with the
\emph{true} $\gu = 2J+1$ and $\gl = 3$ reproduces the NIST oscillator
strengths to five digits:
\[
  f_{\ell u} = 0.05990,\ 0.17974,\ 0.29958
  \quad (J = 0, 1, 2),
\]
in the ratio $1:3:5$ with sum $0.539$ --- the last column of
Table~\ref{tab:weights}.  Had the Table-1 $\gu = 3, 9, 15$ been used in
Eq.~(\ref{eq:Af}) with $\gl = 1$ (the implicit choice behind
Eq.~(\ref{eq:draine8})), the oscillator strengths would be
$f = 0.180, 0.539, 0.899$ (sum $1.62$), a factor of three above the cited
NIST values in every component and in the sum.  The NIST $A$ and $f$ tabulations are therefore mutually
consistent \emph{only} with $\gu = 2J+1$ and $\gl = 3$; the Table-1 $g_u$
column and the $\tau_0$ normalization built on it are the outliers.

%==========================================================================
\section{Independent anchors}
Three further independent references use the $2J+1$ weights and agree with
the standard normalization:
\begin{enumerate}\itemsep2pt
  \item The oscillator-strength form of $\sigma_0$ in
        Eq.~(\ref{eq:sigma0}), derived from first principles for a
        Doppler-broadened line, gives $\tautot^{\rm std} = 87.5$ per
        $10^{14}$\,cm$^{-2}$ at $b = 10$\,km\,s$^{-1}$, identical to the
        $A$-value form.
  \item The p-winds / OH18 exoplanet implementation
        (\code{p\_winds/microphysics.py}) uses the same $f$-values and the
        same $\sigma_0$; that opacity has been validated against
        metastable-helium transit observations.
  \item The LaRT $1.0833\,\mu$m module (\code{LaRT\_v2.00}) adopts the same
        $f = 1:3:5$ ratio and $A = 1.0216\times10^7$\,s$^{-1}$.  A direct
        numerical closure confirms the two independent implementations
        agree: feeding a MoCHII $n(\tSSS)$ grid (the dense H\,\textsc{ii}
        example, center-to-edge $N = 1.71\times10^{15}$\,cm$^{-2}$) into
        LaRT yields \code{tau\_pole} $= 1.294\times10^{3}$ (strongest
        component, center to edge), while MoCHII's blended doublet
        $\tau_0 = 4.12\times10^{3}$ (full diameter) converts to the same
        convention as $4.12\times10^{3} \times \tfrac{1}{2} \times
        f_2/(f_1{+}f_2) = 1.288\times10^{3}$ --- agreement to 0.5\%.
\end{enumerate}
MoCHII's own $\tau_0$ normalization sits in this group.  For thermal
broadening at $T_e = 10^4$\,K ($b_{\rm therm} = 6.45$\,km\,s$^{-1}$,
$v_{\rm turb} = 0$), MoCHII gives a blended line-center opacity of the
$J = 1, 2$ doublet $\sigma_0 = 1.21\times10^{-12}$\,cm$^2$, i.e.\
$\tau_0 = 1$ at $N(\tSSS) = 8.3\times10^{11}$\,cm$^{-2}$, and a doublet-to-
$J{=}0$ optical-depth ratio $(f_1+f_2)/f_0 = 8.00$ --- consistent with the
$f$-value ratios of Table~\ref{tab:weights}.

%==========================================================================
\section{Implications}
The paper's Monte-Carlo profiles are parameterized by $\tautot$, and the
transfer results as functions of $\tautot$ are unaffected --- the issue is
only the mapping from physical conditions to $\tautot$.  For a given
metastable column density, or a given $(Q_{48}, n_3)$, the true optical
depth is one-third of the paper's quoted value.  Table~\ref{tab:tau}
collects the affected coefficients.

\begin{table}[h]
\centering
\caption{Optical-depth coefficients: the paper vs.\ the standard
normalization.  All quantities scale as $(N/10^{14}\,{\rm cm^{-2}})
(10\,{\rm km\,s^{-1}}/b)$ except Eq.~(12), which scales as
$\xi^{1/3} Q_{48}^{1/3} n_3^{4/3}/(1+0.26 n_3)$.}
\label{tab:tau}
\begin{tabular}{lccc}
\toprule
quantity & Draine (2026) & standard / MoCHII & ratio \\
\midrule
single-component coeff.\ (Eq.~9)      & $29.2\,(2J{+}1)$ & $9.72\,(2J{+}1)$ & 3 \\
$\tautot$ coeff.\ (Eq.~11)            & 262   & 87.5 & 3 \\
$\tautot$ from $(Q_{48}, n_3)$ (Eq.~12) & 448 & 149  & 3 \\
\bottomrule
\end{tabular}
\end{table}

The paper's contours and example thresholds should be read accordingly:
the $\tautot$ values quoted for M51 NE-Strip \HII{} ($\sim$180--10$^3$),
M17-B ($\sim$2000), and NGC\,3603 Source E ($\sim$940) are three times the
true values.  The qualitative conclusions nonetheless stand.  Bright \HII{}
regions are optically thick to the $1.0833\,\mu$m triplet under either
normalization --- MoCHII's own $\tau_0$ for representative \HII{} regions
ranges from a classical example (doublet $\tau_0 \approx 34$) to a dense
example ($\tau_0 \sim 4\times10^3$) --- so the predicted broad,
multi-peaked profiles, the $\sim$14\,km\,s$^{-1}$ blueshift, and the
enhancement of dust absorption by resonant trapping all survive.  What
changes is the quantitative placement of the transition thresholds and the
column-density scale inferred from an observed profile: at fixed physical
conditions the resonant-scattering effects set in at three times the
$\tautot$ the paper assigns to those conditions.

Because the printed statistical weights are inconsistent with the NIST
$A$- and $f$-values the paper itself cites, and the standard normalization
is confirmed by three independent implementations, the factor of three
appears to be a transcription of the upper-level degeneracies in Table~1
rather than a difference of convention --- a candidate for an erratum.
MoCHII's $\tau_0$ normalization is unaffected and requires no change.

\vspace{6pt}
\noindent\emph{Reproducibility.}\ All numbers in this memo were recomputed
from the Table-1 $A$-values, wavelengths, and CODATA constants; the check
script is kept in the session scratchpad.

\end{document}
